\title{ Transformation au Rugby }
\author{Marie-Claude, David ; Guillaume, Wuthrich}
\email{marie-claude.david@math.u-psud.fr}
\language{fr}
\precision{2}
\computeanswer{no}
\range{-5..5}

L'essai est marqu� � une distance c du poteau de corner
\integer{c=randint(1..24)}

Le joueur se met � une distance b de la ligne de fond
\integer{b=22+randint(0..8)}

a est la distance au milieu des poteaux
\real{a=sqrt(\b*\b+(35-\c)*(35-\c))}

aa est la distance au milieu du terrain
\real{aa=sqrt((50-\b)*(50-\b)+(35-\c)*(35-\c))}

aaa est la distance � l'entraineur
\real{aaa=sqrt((50-\b)*(50-\b)+(70-\c)*(70-\c))}

\matrix{data=A quelle distance du milieu des poteaux se trouve le ballon, \a
A quelle distance du centre du terrain se trouve le ballon, \aa
L'entra�neur veut donner un conseil � son joueur. Quelle distance s�pare le joueur de l'entra�neur, \aaa}

\text{L=randint(1..rows(\data))}
\text{question=\data[\L;1]}
\real{rep1=\data[\L;2]}
\integer{rep=round(\rep1)}
\statement{<p class="wimscenter">
<img src="\imagedir/rugby.png" /></p><p>
L'essai a �t� marqu� � \(c=\c) m�tres du poteau de corner. <br />
Le joueur d�cide de se placer � \(b=\b) m�tres de la ligne d'essai. <br />
\question ?</p>
<p><i>Donner la r�ponse (en m�tres) sous forme de l'entier le plus proche.</i>}

\answer{La distance en m�tre est d'environ}{\r}{type=numeric }
{option=absolute}
\real{ecart=\reply1-\rep}

\condition{Votre r�ponse est juste � 0.5 m�tre pr�s }{abs(\ecart) <= 0.5}{weight=6}
\condition{Votre r�ponse est juste et respecte la consigne d'approximation }{abs(\ecart) =0}{weight=4}
\feedback{ 1=1}{La bonne r�ponse est \(\rep) m�tres.}

%%%%%%%%%%%%%%%%%%%%%%%%%%%%%%%%%

\condition{Votre r�ponse respecte la consigne d'approximation }{(. notin \reply1) and (/ notin \reply1)}{weight=1}
